\section{The Jacquet functor}
\label{section:Jacquet}

\subsubsection{} In this section we will define a certain category $\Whit(M)_\Conf$ along with functors
  \[\Jac_\Conf : \Whit(G) \to \Whit(M)_\Conf  \qquad  \Poinc_! : \Whit(M)_\Conf \to \DMod(\Bun_M).\]
Furthermore, we will construct a commutative square
\begin{equation}
\label{eqn:Jac_Conf_Poinc_compatibility}
\begin{tikzcd}[
  column sep={9em,between origins},
  row sep={4em,between origins}
]
  \Whit(G) \arrow[d, "\Poinc_!" left] \arrow[r, "\Jac_\Conf"] & \Whit(M)_\Conf \arrow[d, "\Poinc_!"] \\
  \DMod(\Bun_G) \arrow[r, "\CT_{P,!}^-"] & \DMod(\Bun_M) 
\end{tikzcd}
\end{equation}
modulo a certain ULA result whose proof we defer to the next section.

\subsection{Definition of the Jacquet functor}

\subsubsection{Zastava spaces} Write
\[
    \ZZt = \ZZ_{\Nc,\Pt^-} \mon \Bunc_N \underset{\Bun_G}{\times} \Bunt_{P^-}
\]
for the open substack where the generic $N$ and $P^-$ reductions are transverse. We will use the notation $(-)_\ZZt$ to denote restriction to $\ZZt$. We will also use notation such as $\ZZ_{\Bc, \Pt^-}$, $\ZZ_{\Nc, P^-}$, and $\ZZ_{N, P^-}$ to denote the similarly defined stacks with other subgroups and degeneracy conditions.

\subsubsection{Configuration space} Recall that we have a map
  \[\Lambda_G \to \Lambda_{G,P}\]
from the cocharacter lattice of $T$ to the cocharacter lattice of $M^\ab$. We write $\Lambda_{G,P}^\pos$ for the submonoid generated by the images of the simple roots not lying in the Dynkin subdiagram corresponding to $M$. We write
  \[\Conf_{G,P} = \sqcup_{\theta \in \Lambda_{G,P}^\pos} X^\theta\]
for the space of divisors colored by $\alpha_i$, $i \notin \Delta_M$.

\subsubsection{} Let us regard $X$ as a constant curve over $\Conf_{G,P}$ and write $\overset{\circ}{X} \mon X$ for the complement of the tautological divisor.
\begin{itemize}
    \item We define
    \begin{align*}
        \Bunc_{B_M, \Conf}^\vpol &= \Maps^\gen \left(\overset{\circ}{X}, \BB B_M \mon M \bs \left(\overline{M / N_M}\right)/T \right)  \\
        \Bunc_{B_M, \Conf}^\pol &= \Bunc_{B_M, \Conf}^\vpol \underset{\Maps(\overset{\circ}{X}, \BB M)}{\times} \left( \Bun_M \times \Conf \right).
    \end{align*}
    \item Consider the map
        \[\Conf \to \Maps\left(\overset{\circ}{X}, \BB T\right)\]
    given by restricting $\omega^\rhov$ to $\overset{\circ}{X}$. We define
    \begin{align*}
    \Bunc_{N_M, \Conf}^\vpol &= \Bunc_{B_M, \Conf}^\vpol \underset{\Maps(\overset{\circ}{X}, \BB T)}{\times} \Conf  \\
    \Bunc_{N_M, \Conf}^\pol &= \Bunc_{B_M, \Conf}^\pol \underset{\Maps(\overset{\circ}{X}, \BB T)}{\times} \Conf.
    \end{align*}
\end{itemize}
So $\Bunc_{N_M, \Conf}^\vpol$ classifies a point of $\Conf$ along with a degenerate $N_M$-bundle on $\overset{\circ}{X}$. The stack $\Bunc_{N_M, \Conf}^\vpol$ classifies this data along with an extension of the induced $M$-bundle to all of $X$. The reason for introducing $\Bunc_{N_M, \Conf}^\pol$ is that it receives a natural map from $\ZZt$. The following will appear in \S\ref{subsection:Zastava_spaces}.

\begin{construction}
\label{construction:zast_u_lift}
The map $\ZZt \to \Bun_M$ admits a lift
    \[u : \ZZt \to \Bunc_{N_M, \Conf}^\pol.\]
Furthermore, this morphism is proper.
\end{construction}

\subsubsection{} The definition of the full subcategory
    \[
    \Whit(M)_\Conf = \Whit\left(\Bunc_{N_M, \Conf}^\pol\right) \mon \DMod\left(\Bunc_{N_M, \Conf}^\pol\right).
    \]
follows a standard pattern. Write
    \[\left(\Bunc_{N_M, \Conf}^\pol \times \Ran \right)^\good \mon \Bunc_{N_M, \Conf}^\pol \times \Ran\]
for the open subset cut out by the condition that the point of $\Ran$ is disjoint from both the support of the divisor and the degeneracy locus of the $N_M$-reduction. Let
\begin{equation}
\label{eqn:Whit_Conf_Hecke}
\begin{tikzcd}[
  column sep={10em,between origins},
  row sep={2em,between origins}
]
    & \HH_{N_M, \Conf} \arrow[dl] \arrow[dr] \arrow[dd, "\psi"] & \\
    \left(\Bunc_{N_M, \Conf}^\pol \times \Ran \right)^\good & & \left(\Bunc_{N_M, \Conf}^\pol \times \Ran \right)^\good \\
    & \Ga &
\end{tikzcd}
\end{equation}
denote the Hecke groupoid which acts by modifying the $N_M$-bundle, equipped with its natural character. We will say that $\WW$ belongs to $\Whit(M)_\Conf$ if 
    \[(\WW \boxtimes \omega_{\Ran})^\good\]
is equivariant against $\psi^! \cdot \AS$.

\subsubsection{} The functor
    \[
    \oblv : \Whit(M)_\Conf \mon \DMod\left(\Bunc_{N_M, \Conf}^\pol\right)
    \]
admits a continuous right adjoint
    \[
    \Av_* : \DMod\left(\Bunc_{N_M, \Conf}^\pol\right) \to \Whit(M)_\Conf 
    \]
and we define
    \[
    \Poinc_{\Conf,!} = \pi_{\Conf,!} \cdot \oblv
    \qquad
    \coeff_{\Conf,*} = \Av_* \cdot \pi_{\Conf}^!,
    \]
where $\pi_\Conf$ is the projection to $\Bun_M$ displayed in the diagram (\ref{eqn:zastava_diagram}) below.

\subsubsection{} Now we are ready to describe the construction of $\Jac_{\Conf}$. Consider the diagram:
\begin{equation}
\label{eqn:zastava_diagram}
\begin{tikzcd}[
  column sep={4em,between origins},
  row sep={4em,between origins}
]
  & & \ZZt \arrow[dl, hook] \arrow[ddrr, dashed, "u"] & & \\
  & \Bunc_N \! \underset{\Bun_G}{\times} \! \Bunt_{P^-} \arrow[dl, "\aph \widetilde{p}"] \arrow[dr, "\pi"] & & & \\
  \Bunc_N \arrow[dr] & & \Bunt_{P^-} \arrow[dl, "\widetilde{p}"] \arrow[dr, "\qt"] & & \Bunc_{N_M, \Conf}^\pol \arrow[dl, "\pi_\Conf"] \\ 
  & \Bun_G & & \Bun_M &
\end{tikzcd}
\end{equation}
Write $j_{P,!}$ for the $!$-extension of the constant sheaf from $\Bun_{P^-}$ to $\Bunt_{P^-}$. The following assertion is standard, but we will recall its proof at the end of the section.

\begin{proposition}
\label{proposition:zast_vanishing}
Let $i$ denote the closed complement to $\ZZt \mon \Bunc_N \! \underset{\Bun_G}{\times} \! \Bunt_{P^-}$. Then the $!$-pushforward of
  \[i^* \left( \WW \stoximes j_{P,!}^- \right)\]
to $\Bun_M$ is zero.
\end{proposition}

\subsubsection{} This vanishing gives an isomorphism
    \[
    \CT_P^- \cdot \Poinc_! \cdot \WW 
    \simeq (\qt \cdot \pi)_! \cdot \left( \WW \stoximes j_{P,!}^- \right)_\ZZt.
    \]
Since $u$ is proper, the map
    \[
    (\qt \cdot \pi)_! \cdot \left( \WW \stoximes j_{P,!}^- \right)_\ZZt
    \simeq \pi_{\Conf,!} \cdot u_! \cdot \left( \WW \stoximes j_{P,!}^- \right)_\ZZt
    \to \pi_{\Conf,!} \cdot u_* \cdot \left( \WW \shoximes j_{P,!}^- \right)_\ZZt \, [2\delta_G].
    \]
is an isomorphism by the following result, whose proof will occupy \S\ref{section:ULA}.

\begin{lemma}
\label{lem:zast_ULA}
The purity morphism
    \[
    \left( \WW \stoximes j_{P,!}^- \right)_\ZZt 
    \to \left( \WW \shoximes j_{P,!}^- \right)_\ZZt \, [2 \delta_G]
    \]
is an isomorphism.
\end{lemma}

\subsubsection{} Therefore, we define
    \[\Jact_\Conf = u_* \cdot \left( - \stoximes j_{P,!}^- \right)_\ZZ \, [2 \delta_G].\]
and consider the diagram
\[
\begin{tikzcd}
  \Whit(G) \arrow[d] \arrow[r, dashed, "\Jac"] & \Whit(M)_\Conf \arrow[d] \\
  \DMod\left( \Bunc_N \right) \arrow[r, "\Jact"] & \DMod\left( \Bunc_{N_M,\Conf}^\pol \right)
\end{tikzcd}
\]
We define $\Jac_\Conf$ to be the upper dashed arrow, which makes by virtue of the proposition below. It follows that
  \[\CT_{P,!}^- \cdot \Poinc_! \simeq \Poinc_! \cdot \Jac_\Conf.\]

\begin{proposition}
\label{proposition:Jac_Whit_preservation}
The functor $\Jact_\Conf$ takes Whittaker sheaves to Whittaker sheaves.
\end{proposition}

\subsection{Some invariant theory}

\subsubsection{} Construction \ref{construction:zast_u_lift} is based on some elementary facts about the invariant theory of $G$ and its subgroups. Everything we will need is summarized in Proposition \ref{proposition:zast_cartesian}, whose proof can be obtained by a close reading of \cite[\S3]{wang_monoid}.

\subsubsection{Notation} For a classical affine scheme $X = \Spec(A)$ with an action of a group $G$, we write
    \[X \git G = \Spec(A^G)\]
for the naive GIT quotient. Here we are taking invariants in the underived sense. So, for example, we have $\overline{G/N} = G \git N$ in this notation.

\subsubsection{} Recall that the closure of
  \[M \simeq P/U_P \mon G \git U_P\]
is a reductive monoid containing $M$ as its group of units. We will denote it by $M^+$. Similarly, the closure of
  \[M^\ab \simeq P/[P,P] \mon G \git [P,P]\]
is a reductive monoid containing $M^\ab$ as its group of units. We denote it $M^{\ab,+}$. 

\begin{proposition}
\label{proposition:zast_cartesian}
Consider the commutative diagram
\begin{equation}
  \begin{tikzcd}[
  column sep={9em,between origins},
  row sep={4em,between origins}
]
\label{eqn:zast_cartesian}
    U_P^- \bss P^- \overset{M}{\times} P \git N \arrow[r, hook] \arrow[d] & U_P^- \bss G \overset{G}{\times} G \git N \arrow[d] \\
    U_P^- \bss P^- \overset{M}{\times} P \git [P,P] \arrow[r, hook] & U_P^- \bss G \git [P,P]
  \end{tikzcd}
\end{equation}
of affine schemes with $M \times T$ action.
\begin{itemize}
    \item[(i)] This diagram is Cartesian.
    \item[(ii)] The lower row identifies with $M^\ab \mon M^{\ab,+}$.
\end{itemize}
\end{proposition}

\subsubsection{Order of operations} Our convention is that GIT quotient takes precedence over balanced product.

\subsubsection{Proof of Proposition \ref{proposition:zast_cartesian}.(ii)} Actually (ii) is a consequence of a more precise assertion, which we state separately: 

\begin{proposition}
\label{proposition:retract}
The composite
  \[M^{\ab,+} \to G \git [P,P] \to U_P^- \bss G \git [P,P]\]
is an isomorphism.
\end{proposition}

\subsubsection{Proof} Consider the diagram below obtained by functoriality.
\[
\begin{tikzcd}[
  column sep={7em,between origins},
  row sep={3em,between origins}
]
  M^+ \arrow[r, hook] \arrow[d] & G \git U_P \arrow[r] \arrow[d] & U_P^- \bss G \git U_P \arrow[d] \\
  M^+ \git [M,M] \arrow[r, hook] & G \git [P,P] \arrow[r] & U_P \bss G \git [P,P]
\end{tikzcd}
\]
According to \cite[Theorem 3.2.8]{wang_monoid}, the upper row is a retract. Therefore, the lower row is also a retract. In fact, it is a retract of affine schemes, so the lower left arrow is a closed immersion. Therefore, $M^+ \git [M,M]$ contains $M^{\ab,+}$ as a closed subscheme. But $M$ maps dominantly onto both, so they are equal.

\subsubsection{} In order to prove Proposition \ref{proposition:zast_cartesian}.(i), it will be useful to have some generalities on GIT quotients. Let $\UU \mon X$ be integral affine schemes with $G$-action. We will say that the pair $(\UU \mon X, G)$ is \emph{good} if the following square is Cartesian:
  \[
  \begin{tikzcd}[
  column sep={4em,between origins},
  row sep={3em,between origins}
]
    \UU \arrow[r] \arrow[d] & X \arrow[d] \\
    \UU\git G \arrow[r] & X \git G
  \end{tikzcd}
  \]
We will further say that the pair is \emph{quite good} if the lower arrow is a distinguished open affine.

\subsubsection{Example} Let $\UU \mon \AA^2$ be the complement of a line through the origin, and let $\Gm$ act by dilation. Then the pair $(\UU \mon X, \Gm)$ is not good.

\begin{proposition}
\label{proposition:git_qgood}
Let $X$ be an integral affine scheme with $G$-action.
\begin{itemize}
    \item[(i)] For any distinguished open affine $V \mon X \git G$, the pair
      \[\left( V \! \underset{X \git G}{\times} \! X \mon X, G \right)\]
    is quite good.
    \item[(ii)] Suppose that $(\UU \mon X, G)$ is quite good. For any subgroup $H \mon G$, the pair $(\UU \mon X, H)$ is quite good and the square
      \[
      \begin{tikzcd}[
  column sep={4em,between origins},
  row sep={3em,between origins}
]
        \UU \git H \arrow[r, hook] \arrow[d] & X \git H \arrow[d] \\
        \UU \git G \arrow[r, hook] & X \git G
      \end{tikzcd}
      \]
    is Cartesian.
    \item[(iii)] Let $H \mon G$ be a normal subgroup. If $(\UU \mon X, H)$ is good and $(\UU \git H \mon X \git H, G/H)$ is quite good, then $(\UU \mon X, G)$ is quite good.
\end{itemize}
\end{proposition}

\subsubsection{Proof of (i)} Let $\UU$ denote the fiber product indicated by the diagram.
\[
\begin{tikzcd}[
  column sep={3.5em,between origins},
  row sep={3em,between origins}
]
   & \UU \arrow[dl] \arrow[d] \arrow[r, hook] & X \arrow[d] \\
  \UU \git G \arrow[r, dashed] & V \arrow[r, hook] & X \git G
\end{tikzcd}
\]
We need to check that the dashed arrow is an isomorphism. Write $X = \Spec(A)$ and let $f \in A^G$ be a function cutting out the complement of $V$. Then the dashed arrow corresponds to the inclusion of rings
    \[A^G[\tfrac{1}{f}] \mon A[\tfrac{1}{f}]^G.\]
This is an isomorphism.

\subsubsection{Proof of (ii)} Write $V$ for the fiber product appearing in the following Cartesian diagram:
\[
\begin{tikzcd}[
  column sep={4em,between origins},
  row sep={3em,between origins}
]
   & \UU \arrow[dl] \arrow[d] \arrow[r, hook] & X \arrow[d] \\
  \UU \git H \arrow[r, dashed] & V \arrow[r, hook] \arrow[d] & X \git H \arrow[d] \\
   & \UU \git G \arrow[r, hook] & X \git G
\end{tikzcd}
\]
Note that $V \mon X \git H$ is a distinguished open affine because $\UU \git G \mon X \git G$ is. Therefore we may apply (i).

\subsubsection{Proof of (iii)} This is a formal consequence of the pullback lemma.

\subsubsection{Proof of Proposition \ref{proposition:zast_cartesian}.(i)} We give a sequence of reductions.

\subsubsection{First reduction} It suffices to show that (\ref{eqn:zast_cartesian}) is Cartesian after pullback along the $G$-torsor
  \[U_P^- \bss G \times G \git N \to U_P^- \bss G \overset{G}{\times} G \git N.\]
So we need to show that the diagram
\[
\begin{tikzcd}[
  column sep={10em,between origins},
  row sep={4em,between origins}
]
  U_P^- \bss P^- \overset{M}{\times} P \times G \git N \arrow[r, hook] \arrow[d] & U_P^- \bss G \times G \git N \arrow[d] \\
  U_P^- \bss P^- \overset{M}{\times} P \git [P,P] \arrow[r, hook] & U_P^- \bss G \git [P,P]
\end{tikzcd}
\]
is Cartesian. Let us spell out the precise definition of the upper arrow. Consider the pair
\begin{equation}
\label{eqn:git_pair}
    \left(P^- \overset{M}{\times} P \times G \mon G \times G, U_P^- \times G \times [P,P] \right)    
\end{equation}
where the map is
    \[(x, h) \mapsto (x h^{-1}, h).\]
and the group actions are
    \[
    (u, g, p) \cdot (x, h) = (u \cdot x \cdot p, g \cdot h \cdot p)
    \qquad
    (u, g, p) \cdot (h, h') = (u \cdot h \cdot g^{-1}, g \cdot h' \cdot p).
    \]
Taking the GIT quotient by the subgroup $U_P^- \times N \mon U_P^- \times G \times [P,P]$ gives the upper arrow.

\subsubsection{Second reduction} Now Proposition \ref{proposition:git_qgood}.(ii) implies that it suffices to show that (\ref{eqn:git_pair}) is quite good. Note that
    \[
    \left(P^- \overset{M}{\times} P \times G \mon G \times G, G \right)
    \]
is good. Therefore, by Proposition \ref{proposition:git_qgood}(iii), it is enough to prove that
    \[
    \left(P^- \overset{M}{\times} P \mon G, U_P^- \times [P,P] \right)
    \]
is quite good. 

\subsubsection{Third reduction} By Proposition \ref{proposition:retract}, the diagram
\[
\begin{tikzcd}[
  column sep={12em,between origins},
  row sep={5em,between origins}
]
  U_P^- \bss \left( P^- \overset{M}{\times} P \right) \git U_P \arrow[r, hook] \arrow[d] & U_P^- \bss G \git P \arrow[d] \\
  U_P^- \bss \left( P^- \overset{M}{\times} P \right) \git [P,P] \arrow[r, hook] & U_P^- \bss G \git [P,P]
\end{tikzcd}
\]
identifies with
\[
\begin{tikzcd}[
  column sep={5em,between origins},
  row sep={4em,between origins}
]
  M \arrow[r, hook] \arrow[d] & M^+ \arrow[d] \\
  M^\ab \arrow[r, hook] & M^{\ab, +}
\end{tikzcd}
\]
This square is Cartesian by \cite[Appendix C.4.3]{lin_nearby_cycles}. Therefore, by Proposition \ref{proposition:git_qgood}.(iii), it suffices to show that
  \[\left(P^- \overset{M}{\times} P \mon G, U_P^- \times U_P\right)\]
is good. 

\subsubsection{End of the proof} It remains to show that the diagram
\begin{equation}
\label{eqn:git_cartesian}
\begin{tikzcd}[
  column sep={5em,between origins},
  row sep={4em,between origins}
]
  P^- \overset{M}{\times} P \arrow[r, hook] \arrow[d] & G \arrow[d] \\
  M \arrow[r, hook] & M^+ 
\end{tikzcd}
\end{equation}
is Cartesian. For $P = B$ one can verify this by direct calculation \cite[\emph{Proof} of Lemma 3.2.10]{wang_monoid}. The corresponding diagram is displayed on the left.
\[
\begin{tikzcd}[
  column sep={5em,between origins},
  row sep={4em,between origins}
]
  B^- \overset{T}{\times} B \arrow[r, hook] \arrow[d] & G \arrow[d] \\
  T \arrow[r, hook] & N^- \bss G \git N 
\end{tikzcd}
\qquad
\begin{tikzcd}[
  column sep={5em,between origins},
  row sep={4em,between origins}
]
  B^- \overset{T}{\times} B \arrow[r, hook] \arrow[d] & G \arrow[d] \\
  B_M^- \overset{T}{\times} B_M \arrow[r, hook] & M^+ 
\end{tikzcd}
\]
Proposition \ref{proposition:git_qgood}.(ii) now implies that the diagram on the right is also Cartesian. But this diagram will remain Cartesian if we replace the left vertical arrow with its $M$-orbit. The result of this operation is (\ref{eqn:git_cartesian}).

\subsection{Zastava spaces}
\label{subsection:Zastava_spaces}

\subsubsection{} Now we will obtain a precise description of $\ZZt$ in local terms. In order to state the result, we will need to introduce some more stacks over $\Conf_{G,P}$. 

\subsubsection{} As before, let us view $X$ as a constant curve over $\Conf$. Write
    \[\Disco \mon \Disc \mon X\]
for the punctured formal disc inside the formal disc around the support of the tautological divisor. For a map $\YY_1 \to \YY_2$, we use the notation
    \[
    \Maps^\gen(\Disc, \YY_1 \to \YY_2) = \Maps(\Disc, \YY_2) \underset{\Maps(\Disco, \YY_2)}{\times} \Maps(\Disco, \YY_1).
    \]

\subsubsection{Notation}
\begin{itemize}
    \item We define
        \begin{align*}
            \Gr_{G,\Conf} &= \Maps^\gen \left( \Disc, \pt \to \BB G \right)   \\
            \Grc_{B,\Conf} &= \Maps^\gen \left( \Disc, \pt \to T \bs (N \bss G) / G \right)   \\
            \Grt_{P^-,\Conf} &= \Maps^\gen \left( \Disc, \pt \to G \bs (G \git U_P^- ) / M \right)
        \end{align*}
    \item Note that there is a tautological section
        \begin{equation}
            \label{eqn:tautological_section}
            \Conf \to \Maps^\gen(\Disc, M^\ab \bs M^{\ab, +}).
        \end{equation}
    The map
        \[N \bss G \overset{G}{\times} G \git U_P^- \to M^{\ab, +}\]
    induces a morphism
        \[\Grc_{B,\Conf} \underset{\Gr_{G,\Conf}}{\times} \Grt_{P^-, \Conf} \to \Maps^\gen(\Disc, M^\ab \bs M^{\ab, +})\]
    and we define $\ZZ_{\Bc, \Pt, \Conf}$ to be its base change along (\ref{eqn:tautological_section}).
\end{itemize}
\subsubsection{Variants}
\begin{itemize}
    \item Consider the map
        \[
        \Conf \to \Maps\left(\overset{\circ}{X}, \BB T \right)
        \]
    given by restricting $\omega^\rhov$ to $\overset{\circ}{X}$. We define
        \[
        \Bunc_{N_M, \Conf}^\vpol = \Bunc_{B_M, \Conf}^\vpol \underset{\Maps(\overset{\circ}{X}, \BB T)}{\times} \Conf.
        \]
    \item Consider the map
        \[
        \Conf \to \Maps(\Disc, \BB T)
        \]
    given by restricting $\omega^\rhov$ to $\DD$. We define
        \begin{align*}
            \Grc_{N,\Conf} &= \Grc_{B,\Conf} \underset{\Maps(\Disc, \BB T)}{\times} \{\omega_\Disc^\rhov\} \\
            \ZZ_{\Nc, \Pt^-, \Conf} &= \ZZ_{\Bc, \Pt^-, \Conf} \underset{\Maps(\Disc, \BB T)}{\times} \{\omega_\Disc^\rhov\}
        \end{align*}
\end{itemize}

\begin{proposition}
\label{proposition:zast_loc_formula}
Define $\ZZ_{\Bc, \Pt^-}$ by replacing $N$ with $B$ in the definition of $\ZZt = \ZZ_{\Nc, \Pt^-}$. There are canonical isomorphisms
    \begin{align*}
        \ZZ_{\Bc, \Pt^-} & \simeq \Bunc_{B_M, \Conf}^{\vpol} \timest \ZZ_{\Bc, \Pt^-, \Conf} \\
        \ZZ_{\Nc, \Pt^-} & \simeq \Bunc_{N_M, \Conf}^{\vpol} \timest \ZZ_{\Nc, \Pt^-, \Conf}.
    \end{align*}
\end{proposition}

\subsubsection{} Let us explain the meaning of these twisted products, which \emph{prima facie} involve fiber products over $\BB \LL_\Conf B_M$ and $\BB \LL_\Conf N_M^\omega$. To avoid the problem of taking quotients by ind-infinite dimensional group schemes, we use
    \[\Maps(\Disco, \BB B_M)\]
as a substitute for $\BB \LL_\Conf B_M$. Then the twisted product involving $B$ is defined to be the pullback of 
\begin{equation}
\label{eqn:Zast_descent}
    \Maps^\gen\left(\Disc, \BB B_M \mon T \bs \left( N \bss G \overset{G}{\times} G \git U_P^- \right) / M \right)
\end{equation}
along the natural map
    \[
    \alpha : \Bunc_{B_M, \Conf}^\vpol \to \Maps(\Disco, \BB B_M).
    \]
The second twisted product involving $N$ is defined similarly.

\subsubsection{Sanity check} The expression in (\ref{eqn:Zast_descent}) is a substitute for
    \[\LL_\Conf B_M \bs \ZZ_{\Bc, \Pt^-, \Conf}\]
in the sense that its base change along
    \[\beta : \Conf \to \Maps(\Disco, \BB B_M)\]
identifies with $\ZZ_{\Bc, \Pt^-, \Conf}$. Since $\alpha$ Zariski locally factors through $\beta$, the twisted product is Zariski locally on $\Bunc_{B_M,\Conf}^\vpol$ a product
    \[\Bunc_{B_M,\Conf}^\vpol \underset{\Conf}{\times} \ZZ_{\Bc, \Pt^-, \Conf}\]
over the configuration space.

\subsubsection{Completion of Construction \ref{construction:zast_u_lift}} Using the identification of Proposition \ref{proposition:zast_loc_formula}, we define the morphism $u$ to be
    \[
    \Bunc_{N_M, \Conf}^{\vpol} \timest \ZZ_{\Nc, \Pt^-, \Conf}
    \to \Bunc_{N_M, \Conf}^{\vpol} \timest \Gr_{M,\Conf}
    \simeq \Bunc_{N_M, \Conf}^\pol.
    \]
It is proper because $\ZZ_{\Nc, \Pt^-, \Conf} \to \Conf$ is proper.

\subsubsection{Proof of Proposition \ref{proposition:zast_loc_formula}} Let us rewrite (\ref{eqn:zast_cartesian}) as
\[
  \begin{tikzcd}[
  column sep={7em,between origins},
  row sep={4em,between origins}
]
    M \git N_M \arrow[r, hook] \arrow[d] & U_P^- \bss G \overset{G}{\times} G \git N \arrow[d] \\
    M^\ab \arrow[r, hook] & M^{\ab,+}
  \end{tikzcd}
\]
The quotient of the lower row by $M \times T$ identifies with
    \[\BB [M,M] \times \BB T \mon \BB [M,M] \times M^\ab \bs M^{\ab,+} \times \BB T.\]
Therefore, we obtain a Cartesian diagram
\[
  \begin{tikzcd}[
  column sep={12em,between origins},
  row sep={4em,between origins}
]
    M \bs \left( M \git N_M \right) / T \arrow[r, hook] \arrow[d] & M \bs \left( U_P^- \bss G \overset{G}{\times} G \git N \right) / T \arrow[d] \\
    \pt \arrow[r, hook] & M^\ab \bs M^{\ab,+}
  \end{tikzcd}
\]
Now Proposition \ref{proposition:zast_loc_formula} is a formal consequence of Beauville--Laszlo gluing.

\subsection{The Whittaker condition}
\label{subsection:Jacquet_Whittaker}

\subsubsection{Proof of Proposition \ref{proposition:zast_vanishing}} We follow the usual pattern of passing to strata and applying character orthogonality.

\subsubsection{First reduction} The object $\WW \stoximes j_{P,!}$ in question is $!$ extended from the locus
  \[\Bun_N \underset{\Bun_G}{\times} \Bun_{P^-}.\]
Therefore, it suffices to show that for each locally closed substack
  \[i_w : \ZZ_{N,P^-}^w = \left( \Bun_N \underset{\Bun_G}{\times} \Bun_{P^-} \right)^w \mon \Bun_N \underset{\Bun_G}{\times} \Bun_{P^-}\]
with $w \neq 1$, the object
  \[i_w^* \left( \WW \stoximes k_{\Bun_{P^-}} \right) \simeq \aph p^* \cdot \WW \]
is annihilated by $!$-pushforward to $\Bun_M$. 

\subsubsection{Second reduction} For any finite set $x \mon X$ we may consider the open substack
  \[\ZZ_{N,P^-}^{w, \good\,x} \mon \ZZ_{N,P^-}^w\]
where the $N$ and $P^-$ reductions are in relative position $w$ at $x$. Since $\ZZ_{N,P^-}^w$ is a colimit of open substacks of this form, it suffices to show that the $*$-pullback of $\WW$ to each $\ZZ_{N,P^-}^{w, \good\,x}$ is annihilated by $!$-pushforward to $\Bun_M$. 

\subsubsection{End of the proof} Now consider the groupoid acting on $\ZZ_{N,P^-}^{w, \good \, x}$ by $N^w \cap P^-$-Hecke modifications at $x$. The unipotent subgroupoid corresponding to
    \[N^w \cap U_P^- \simeq \ker\left( N^w \cap P^- \to M \right)\]
acts along the fibers of the projection to $\Bun_M$. Now the $*$-pullback of $\WW$ is equivariant for the character associated to the pullback of $\psi$ along
  \[N^w \cap U_P^- \mon N^w \simeq N \to \Ga,\]
which is nontrivial because $w \neq 1$.

\subsubsection{Proof of Proposition \ref{proposition:Jac_Whit_preservation}} This claim is verified by doing some bookkeeping with Hecke groupoids.

\subsubsection{Step one} We form an extension of the diagram (\ref{eqn:Whit_Conf_Hecke}) appearing in the definition of Whittaker condition for $\Whit(M)_\Conf$.
\[
\begin{tikzcd}[
  column sep={10em,between origins},
  row sep={2em,between origins}
]
    & \HH_\ZZt \arrow[dl] \arrow[dr] \arrow[dd] & \\
    \left(\ZZt \times \Ran \right)^\good \arrow[dd] & & \left(\ZZt \times \Ran \right)^\good \arrow[dd] \\
    & \HH_{N_M, \Conf} \arrow[dl] \arrow[dr] \arrow[dd, "\psi"] & \\
    \left(\Bunc_{N_M, \Conf}^\pol \times \Ran \right)^\good & & \left(\Bunc_{N_M, \Conf}^\pol \times \Ran \right)^\good \\
    & \Ga &
\end{tikzcd}
\]
Here are the definitions:
\begin{itemize}
    \item We write $(\ZZt \times \Ran)^\good$ for the preimage of
        \[
        \left(\Bunc_{N_M, \Conf}^\pol \times \Ran \right)^\good \mon \Bunc_{N_M, \Conf}^\pol \times \Ran
        \]
    along the projection
        \[u \times \id : \ZZt \times \Ran \to \Bunc_{N_M,\Conf}^\pol \times \Ran.\]
    \item The fiber products of the two faces are naturally identified. We write $\HH_\ZZt$ for this common value; it is naturally a groupoid acting on $(\ZZt \times \Ran)^\good$.
\end{itemize}
Therefore, it remains to show that
    \[\left(\WW \shoximes j_{P,!}^-\right)_\ZZt\]
is $(\HH_\ZZt, \psi)$-equivariant, since $u_*$ sends such objects into $\Whit(M)_\Conf$.

\subsubsection{Step two} The $!$-pullback of $\WW$ along $\ZZt \to \Bunc_N$ is $(\HH_\ZZt, \psi)$-equivariant. Indeed, we have a commutative diagram
\[
\begin{tikzcd}[
  column sep={8em,between origins},
  row sep={2em,between origins}
]
    & \HH_\ZZt \arrow[dl] \arrow[dr] \arrow[dd, dashed] & \\
    \left(\ZZt \times \Ran \right)^\good \arrow[dd] & & \left(\ZZt \times \Ran \right)^\good \arrow[dd] \\
    & \HH_N \arrow[dl] \arrow[dr] \arrow[dd, "\psi"] & \\
    \left(\Bunc_N \times \Ran \right)^\good & & \left(\Bunc_N \times \Ran \right)^\good \\
    & \Ga &
\end{tikzcd}
\]
in which the two faces are not Cartesian, and the pullback of the Whittaker character from $\HH_N$ agrees with the pullback of the Whittaker character from $\HH_{N_M,\Conf}$.

\subsubsection{Step three} The $!$-pullback of $j_{P,!}$ along $\ZZt \to \Bunt_{P^-}$ is $\HH_\ZZt$-equivariant. Indeed, we have a commutative diagram
\[
\begin{tikzcd}[
  column sep={8em,between origins},
  row sep={2em,between origins}
]
    & \HH_\ZZt \arrow[dl] \arrow[dr] \arrow[dd, dashed] & \\
    \left(\ZZt \times \Ran \right)^\good \arrow[dd] & & \left(\ZZt \times \Ran \right)^\good \arrow[dd] \\
    & \HH_{P^-} \arrow[dl] \arrow[dr] & \\
    \left(\Bunt_{P^-} \times \Ran \right)^\good & & \left(\Bunt_{P^-} \times \Ran \right)^\good
\end{tikzcd}
\]
in which the two faces are not Cartesian, and $j_{P,!}^-$ is equivariant for $\HH_{P^-}$.
